%% appE-further-reading \dash{} Co dalej
%
% NIC TUTAJ NIE JEST PISANE Z PAMIĘCI i to ograniczenie ukształtowało ten
% dodatek bardziej niż cokolwiek innego. Sekcje E.1 i E.2 to przegląd
% konkurencji z notes/02-grounding-and-traps.md sekcja 1, przeprowadzony przed
% zaplanowaniem książki. Podaje źródła twierdzeń o świecie; jego sekcja Sources
% nazywa recenzję dla jednego z dwóch niesionych tu cytatów, a dla
% "encyklopedycznej pozycji referencyjnej" Murphy'ego nie nazywa żadnej, więc
% dodatek niczego nie obiecuje. Sekcja E.3 to notes/01-curriculum.md sekcja 13.
%
% Sekcja E.4 podaje NAZWY, a nazwa jest tam werdyktem tylko tam, gdzie E.1 już
% go wydała -- Strang, Deisenroth/Faisal/Ong, Goodfellow/Bengio/Courville,
% Murphy i Bishop powtarzają się w obu. TO POKRYCIE JEST POWODEM BRZMIENIA
% RAMKI OSTRZEŻENIA: mówiła kiedyś wprost, że poniższa lista nie została
% przeczytana i porównana, co jej własne nazwy fałszowały. Podawaj regułę,
% nigdy zliczenia. Dla reszty repozytorium notuje tytuł i nic więcej, więc
% dopisanie charakterystyki byłoby jej wymyśleniem -- a to jest dokładnie ta
% klasa, przed którą książka przestrzega przez czterdzieści siedem programów.
%
% Każde twierdzenie o PROGRAMIE zostało tu sprawdzone w napisanym programie,
% zanim trafiło na stronę -- to dyscyplina wypracowana w passie Dodatku D.
% Nic tego nie bramkuje; nic nie może.
%
%
% SEKCJA E.5 TO SKOROWIDZ TRAS DLA RAMEK RYGORU i jest bramkowana:
% tools/check_structure.py --rigour przewraca build, gdy ramka odkłada wynik
% poza książkę, a żaden wiersz tutaj mu nie odpowiada. Ramka odkładająca
% WEWNĄTRZ książki cytuje program, do którego odkłada, więc ramka bez
% \ref{prog:...} jest z konstrukcji skierowana na zewnątrz -- to kryterium jest
% dokładne, jest stabilne między wydaniami i to je czyta ten test. Czego NIE
% umie przeczytać, to czy nazwany cel zawiera dowód; pięć wierszy mówi "nie
% zbadane" właśnie dlatego, że nikt tej lektury nie odbył.
%
% WCIĄŻ NIEZAŁATWIONE: nic w książce nie odwołuje się do tego dodatku. Nigdzie
% w programs/ nie ma \ref{app:E}, więc czytelnik stojący przy ramce rygoru nie
% ma jak się dowiedzieć, że sekcja E.5 istnieje. Domknięcie tego wymaga edycji
% każdego odkładającego programu i jest passem samym w sobie.
%
% A OŚWIADCZENIE O ZAKRESIE MA DWA ŹRÓDŁA, nie jedno. Dwie pierwsze tabele
% sekcji E.3 to notes/01-curriculum.md sekcja 13; trzecia tabela i cztery
% akapity "najmniej pewne" to notes/02-grounding-and-traps.md sekcja 4, czyli
% druga i większa lista wyłączeń, której pierwsza wersja tego dodatku nigdy nie
% otworzyła. Rozszerzając oświadczenie o zakresie, otwórz obie.
%
% ŻADNYCH ADRESÓW WWW. Treść tej książki nie drukuje żadnego, a wydrukowany
% adres psuje się szybciej niż tytuł.

\chapter{Co dalej}
\label{app:E}

\noindent
Ten dodatek składa się z dwóch połówek, a różnica między nimi jest powiedziana,
a nie ukryta.

Sekcje~\ref{sec:E-survey} i~\ref{sec:E-positions} podają werdykt o każdej
pozycji, bo te prace zostały przeczytane i porównane, zanim książka powstała:
przegląd, z którego pochodzą, istnieje po to, by odpowiedzieć na dwa pytania
\dash{} \emph{co już istnieje}, oraz \emph{co każda z tych prac robi dobrze,
a co źle dla tego czytelnika?} Tam, gdzie cytat nie jest tu przypisany do
nazwanej recenzji, jest tak dlatego, że zapis, z którego pochodzi, żadnej nie
nazywa.
Sekcja~\ref{sec:E-not} mówi, czego ta książka nie uczy, dlaczego i gdzie dany
temat jest zrobiony porządnie.

Sekcja~\ref{sec:E-depth} to lista nazw i mówi, które z nich mają werdykt.
Część to prace, które Sekcja~\ref{sec:E-survey} już oceniła, i tamta ocena
zostaje; reszty ta książka nie sprawdziła tak samo, a jednozdaniowa
charakterystyka wymyślona po to, by wypełnić lukę, byłaby dokładnie tym,
o czym czterdzieści siedem programów prosi cię, byś tego nie przyjmował.

Nie ma tu adresów WWW. Tytuł i nazwisko przeżyją każdy URL.

\section{Książki, wobec których mierzona jest ta}
\label{sec:E-survey}
\index{further reading!competitive survey}

Uporządkowane według tego, jak blisko każda praca stoi czytelnika, dla którego
ta książka powstała: pracującego inżyniera, sprawnego w kodzie, któremu
matematyka ze szkoły wyblakła, a który dziś wdraża modele.

Każdy wpis mówi, co dana praca robi dobrze, co robi źle \emph{dla tego
czytelnika}, jaki próg zakłada przed pierwszą stroną i czy jest pierwszą
lekturą, czy odnośnikiem. \textbf{Żaden nie mówi, ile czasu zajmuje}, i to
pominięcie jest celowe, a nie przeoczone: nikt w tym projekcie nie zmierzył
żadnej lektury, a czas wymyślony po to, żeby wypełnić kolumnę, byłby pierwszym
niepopartym twierdzeniem w tej książce. To próg i skala rozstrzygają, a one są
faktami.

\medskip
\textbf{Deisenroth, Faisal i Ong, \emph{Mathematics for Machine Learning}.}
Najbliższy konkurent i nazywa tę samą lukę co ta książka \dash{} autorzy piszą,
że dystans między matematyką szkolną a poziomem potrzebnym do czytania
podręcznika uczenia maszynowego jest dla wielu osób za duży.
\textbf{Robi dobrze:} najlepiej zorganizowana mapa terenu, darmowy PDF, spójna
narracja od przestrzeni wektorowych do czterech opracowanych metod i najlepsze
krótkie ujęcie rozkładów macierzy, jakie jest dostępne.
\textbf{Robi źle, dla tego czytelnika:} nazywa lukę szkoła\dash{}uczenie
maszynowe swoim celem, a potem zaczyna powyżej niej \dash{} rozdział 2 otwiera
się układami równań liniowych, a nic w tej książce nie uczy przekształcania
nierówności ani przesuwania wskaźnika sumowania \dash{} progiem, jaki w
praktyce zakłada, jest sprawna algebra i trochę rachunku różniczkowego.
Cztery docelowe metody są
sprzed uczenia głębokiego, nie ma ani jednej strony o nadmiarze czy
uwarunkowaniu, a większość zadań nie ma rozwiązania.
\textbf{Werdykt:} najlepsza pozycja, na którą można wskazać, i mimo przedmowy
nie jest drogą od zera.

\medskip
\textbf{Goodfellow, Bengio i Courville, \emph{Deep Learning}, część~I.}
\textbf{Robi dobrze:} rozdział 4, \emph{Numerical Computation}, to
najwartościowszy rozdział w całym tym przeglądzie i jedyny, którego
nikt nie naśladuje \dash{} nadmiar i niedomiar, stabilizacja log-sum-exp,
uwarunkowanie, jakobian i hesjan. To jedyna praca tutaj, która traktuje maszynę
jako obiekt o skończonej precyzji. \textbf{Robi źle:} część~I jest
przypomnieniem, a nie nauczaniem \dash{} cztery rozdziały w miejsce dwóch lat
studiów, co jest przypomnieniem dla kogoś, kto ten materiał spotkał, a nie
nauczaniem dla kogoś, kto go nie spotkał, bez zadań i bez przykładów
rozwiązanych krok po kroku \dash{} i jest
zamrożona w 2016 roku, więc nie ma transformera ani współczesnej praktyki
mieszanej precyzji. \textbf{Werdykt:} ten rozdział jest przodkiem części~II tej
książki, a ta książka schodzi o rząd wielkości głębiej i z liczbami.

\medskip
\textbf{Strang, \emph{Linear Algebra and Learning from Data}.}
\textbf{Robi dobrze:} Strang jest najlepszym żyjącym wykładowcą algebry
liniowej, a to jest miejsce, w którym łączy cztery podprzestrzenie, SVD i
przybliżenie niskiego rzędu z danymi. Ustawienie SVD w środku ciężkości jest
dokładnie słuszne.
\textbf{Robi źle:} recenzenci nie owijają w bawełnę \dash{} recenzja MAA pisze,
że autor próbował wtłoczyć zdecydowanie za dużo materiału w 432 strony, więc
każdy temat dostaje bardzo szybkie ujęcie, a wyniki bywają podane bez dowodu i
bez cytowania. Zakłada też ukończony pierwszy kurs; to kurs drugi.
\textbf{Werdykt:} źródło ujęcia, a nie wzór konstrukcji.

\medskip
\textbf{Boyd i Vandenberghe, \emph{Introduction to Applied Linear Algebra}.}
\textbf{Robi dobrze:} najlepiej zaprojektowany podręcznik stosowanej algebry
liniowej w tym przeglądzie i temperamentem najbliższy tej książce. Jest
agresywnie wąski \dash{}
zbudowany na liniowej niezależności i faktoryzacji QR i prawie na niczym innym
\dash{} i dowodzi, że książka może pominąć większość swojego przedmiotu i być
przez to lepsza. Darmowy PDF, z materiałami w Pythonie i Julii.
\textbf{Robi źle, dla tego czytelnika:} z założenia pomija wartości własne, SVD
i wyznaczniki, czyli dokładnie te części, których potrzeba do uwarunkowania,
uwagi i kowariancji, a i tak zakłada, że algebra przychodzi ci bez wysiłku.
\textbf{Werdykt:} dyscyplina zakresu do naśladowania. Jego pominięcia nie są
pominięciami tej książki.

\medskip
\textbf{Boyd i Vandenberghe, \emph{Convex Optimization}.}
\textbf{Robi dobrze:} pozycja referencyjna. Wypukłość, dualność, a przede
wszystkim problem rozpoznania \dash{} \emph{czy mój problem jest wypukły?}
\dash{} to trwale przydatne słownictwo. Darmowy PDF, obszerne zadania.
\textbf{Robi źle, dla tego czytelnika:} sama książka mówi, że chce dobrej
algebry liniowej na wejściu, a uczenie głębokie nie należy do jej klasy
problemów. Cele niewypukłe, stochastyczne i przeparametryzowane leżą poza tą
teorią, a dowody tempa zbieżności się nie przenoszą.
\textbf{Werdykt:} weź słownictwo do Programów~\ref{prog:P19}
i~\ref{prog:P22}; nie bierz programu nauczania.

\medskip
\textbf{Bishop, \emph{Pattern Recognition and Machine Learning}, oraz Bishop i
Bishop, \emph{Deep Learning: Foundations and Concepts}.}
\textbf{Robi dobrze:} organizacja, rysunki i materiał o modelach
probabilistycznych należą do najlepiej napisanych gdziekolwiek, a tom z 2024
roku przenosi to samo ujęcie na transformery i dyfuzję. \textbf{Robi źle:} obie
są bardziej zaawansowane niż większość wprowadzeń i zakładają rachunek wielu
zmiennych, algebrę liniową i rachunek prawdopodobieństwa jako narzędzia
robocze; matematyka jest tam zawsze w służbie modelu i nigdy nie jest uczona od
podłogi. \textbf{Werdykt:} cel, a nie droga \dash{} a warunkiem powodzenia tej
książki jest to, że możesz ją teraz otworzyć.

\medskip
\textbf{Murphy, \emph{Probabilistic Machine Learning} \dash{} wprowadzenie
z 2022 roku i tom zaawansowany z 2023.}
\textbf{Robi dobrze:} najpełniejsze i najbardziej aktualne pojedyncze ujęcie,
spięte modelowaniem probabilistycznym, z dobrą sekcją o tle matematycznym i
otwartym kodem. \textbf{Robi źle:} poziom magisterski wprost z własnej
deklaracji, z założeniem rachunku różniczkowego, statystyki i algebry liniowej;
opis recenzenta, że to encyklopedyczna pozycja referencyjna, jest zarazem
pochwałą i problemem, przy ponad tysiącu stron w obu tomach razem.
\textbf{Werdykt:} pozycja referencyjna, do której kierować się tematami, kiedy
już umiesz ją czytać.

\medskip
\textbf{Blum, Hopcroft i Kannan, \emph{Foundations of Data Science}.}
\textbf{Robi dobrze:} jedyna książka na tej liście, która poważnie traktuje to,
co inżynierowie mylą najbardziej niezawodnie \dash{} zachowanie danych w
wysokim wymiarze. Koncentracja miary, rzutowanie losowe, SVD, błądzenia losowe.
\textbf{Robi źle, dla tego czytelnika:} to kurs teoretyczny, prowadzony
dowodami i adresowany do studentów teorii informatyki, a jego postęp nie
przypomina potrzeb praktyka.
\textbf{Werdykt:} źródło treści Programu~\ref{prog:P05}, którą ta książka
przekłada z postaci twierdzenie\dash{}dowód na coś, na czym można działać.

\section{Dwa stanowiska i jedna metoda}
\label{sec:E-positions}
\index{further reading!positions}

\medskip
\textbf{3Blue1Brown.} \emph{Essence of Linear Algebra} i \emph{Essence of
Calculus} robią coś, czego nie robią tutejsze książki: sprawiają, że
geometryczna treść wyznacznika albo pochodnej staje się widoczna w kilka
sekund.

A najmocniejszy argument za czytaniem książki takiej jak ta stawia ich autor.
Sanderson mówił publicznie, że obawia się, iż jego praca jest intelektualnym
cukierkiem dającym poczucie rozumienia, i że jeśli widz nie zaangażuje się
czynnie \dash{} nie wymyśli tego od nowa, nie porozwiązuje zadań \dash{} to mu
nie zostanie; rozciąga tę uwagę na wykłady Feynmana, mówiąc, że to, jak
przyjemnie ogląda się dany materiał, nie koreluje z uczeniem się na dłuższą
metę.

To ustępstwo osoby, która ma na nim najwięcej do stracenia, i jest to teza tej
książki w jednym zdaniu. Właściwa relacja jest komplementarna: obejrzyj film dla
obrazu, przerób program dla zapamiętania.

\medskip
\textbf{fast.ai i stanowisko, że matematyki potrzeba mniej, niż ci się
wydaje.} Ma rację co do \emph{wejścia} i ta książka przyznaje to bez
zastrzeżeń: twierdzenie, że dwa lata matematyki muszą poprzedzić wytrenowanie
klasyfikatora, odcinało ludzi od tej dziedziny, nie dało nic i jest empirycznie
fałszywe.

To stanowisko odpowiada na pytanie, \emph{czy mogę coś wdrożyć?} Nie odpowiada,
dlaczego strata poszła w \code{nan} na kroku cztery tysiące i tylko w połowicznej
precyzji; czy trzypunktowy zysk w ewaluacji jest prawdziwy, czy to jedno
losowanie z zaszumionego rozkładu; ani czy ustawienie akumulacji gradientu po
cichu zmieniło funkcję celu. Te pytania przychodzą \emph{po} wdrożeniu i są
pytaniami matematycznymi.

Zatem: nie ucz się najpierw matematyki. Ta książka jest tym, po co sięgasz
jako drugie, kiedy coś nie działa, a ty nie widzisz dlaczego.

\medskip
\textbf{Stroud, \emph{Engineering Mathematics} \dash{} metoda, a nie
konkurent.} Ramka, odpowiedź na następnej stronie, Quiz, Podsumowanie i lista
\emph{Czy potrafisz?} są jego. Jego własne ograniczenia są zobowiązaniami tej
książki i są nazwane we wstępie, a nie ukryte:
brak ścisłości; algebra liniowa na poziomie przepisów, bez przestrzeni
wektorowych i bez rozkładów; statystyka opisowa bez wnioskowania; brak
matematyki dyskretnej; brak optymalizacji w pierwszym tomie; oraz \dash{}
zarzut stawiany najczęściej \dash{} bezużyteczność jako pozycji referencyjnej po
przeczytaniu. Część~III odpowiada na algebrę liniową, część~IV na matematykę
dyskretną i ścisłość, część~VI na optymalizację, a część~VII na wnioskowanie;
Dodatki~B, C i D istnieją dla ostatniego zarzutu.

\section{Czego ta książka nie uczy i gdzie iść zamiast tego}
\label{sec:E-not}
\index{further reading!exclusions}

Reguła domu brzmi: książka może powiedzieć, że dany temat nie jest wart swojego
kosztu dla tego czytelnika. Nie może udawać, że temat nie istnieje.

Najpierw ścisłe fundamenty pod tym, czego ta książka uczy \dash{} miejsca, z
których bierze wynik i go nie dowodzi.

\begin{center}
\begin{tabularx}{\linewidth}{@{}XXl@{}}
\toprule
\textbf{Czego nie uczy} & \textbf{Dlaczego} & \textbf{Gdzie iść} \\
\midrule
Rachunku prawdopodobieństwa w ujęciu teorii miary & Nic, co ta książka spłaca,
  nie potrzebuje sigma-ciała, a ta maszyneria kosztuje semestr & Williams;
  Durrett \\
Analizy i dowodzenia na poziomie epsilon--delta & Program~\ref{prog:P14} uczy
  \emph{czytać} twierdzenie i celowo nie uczy go pisać & Abbott; Tao \\
Równań różniczkowych cząstkowych & Czytelnik spotyka je w pracy z fizyką w
  modelu, a płytkie ujęcie tego nie obsłuży & Strauss; Evans \\
Stochastycznych równań różniczkowych, a więc i pełnej matematyki modeli
  dyfuzyjnych & Książka daje rozkład Gaussa i fakt, że wariancje niezależnych
  wielkości się dodają, co jest całym łańcuchem Gaussów w czasie dyskretnym, i
  nigdy tego do niego nie stosuje & Särkkä i Solin \\
Analizy zespolonej i transformaty Fouriera, o której ta książka nie wspomina &
  Program~\ref{prog:F08} daje okrąg jednostkowy; twierdzenia o splocie to
  osobna książka & Bracewell \\
\bottomrule
\end{tabularx}
\end{center}

\medskip
A dalej specjalizacje, które zaczynają się tam, gdzie kończy się głębokość
robocza. Każda z nich to osobna dziedzina, a książka zatrzymuje się w miejscu,
w którym inżynier przestaje potrzebować się o nią spierać.

\begin{center}
\begin{tabularx}{\linewidth}{@{}XXl@{}}
\toprule
\textbf{Czego nie uczy} & \textbf{Dlaczego} & \textbf{Gdzie iść} \\
\midrule
Numerycznej algebry liniowej na poziomie implementacji & Program~\ref{prog:P09}
  mówi, czemu nie odwracać, a Program~\ref{prog:P11}, czemu równania normalne
  podnoszą wskaźnik uwarunkowania do kwadratu; żaden nie uczy implementować
  faktoryzacji QR & Trefethen i Bau; Higham \\
Teorii optymalizacji wypukłej i pełnej dualności & Programy~\ref{prog:P19}
  i~\ref{prog:P22} dają część roboczą, a nie ma powodu parafrazować jednej z
  najlepiej napisanych książek matematycznych & Boyd i Vandenberghe \\
Teorii uczenia statystycznego: wymiaru VC, ograniczeń na uogólnianie &
  Program~\ref{prog:P14} nazywa je twierdzeniami do czytania uważnie; te
  ograniczenia są w realnej skali prawie zawsze puste & Shalev-Shwartz i
  Ben-David \\
Teorii uczenia ze wzmocnieniem i wyników o zbieżności & Programy~\ref{prog:P21}
  i~\ref{prog:P22} dają estymatory gradientu i ograniczenie KL, czyli to, czym
  inżynier faktycznie manipuluje & Sutton i Barto \\
Teorii kategorii, geometrii różniczkowej, rozmaitości & Czasem przywoływane w
  pracach o interpretowalności; nie są nośne dla tej roboty, a na teorię
  kategorii nie przejrzano tu żadnego źródła & Lee, na geometrię \\
\bottomrule
\end{tabularx}
\end{center}

\medskip
A na koniec klasyczny program, z którego ta książka rezygnuje. Dwie tabele
powyżej to tematy, z których \emph{pożycza}, więc każdy wskazuje książkę. Tu są
techniki, które kurs matematyki wyćwiczyłby, a ich celem nie jest żadna książka
\dash{} jest nim program tej książki i procedura, którą już masz zainstalowaną.
Na tym polega cały argument za ich wycięciem: nie na tym, że są nieciekawe, ale
na tym, że robotę, do której wymyślono je jako robotę ręczną, robi dziś
biblioteka, a umiejętnością, która przeżywa, jest wiedza, które wywołanie i
dlaczego.

\begin{center}
\begin{tabularx}{\linewidth}{@{}XX@{}}
\toprule
\textbf{Niećwiczone} & \textbf{Co robi tę robotę zamiast tego} \\
\midrule
Techniki całkowania: przez części, przez podstawienie, ułamki proste & Nic w
  pętli treningowej nie całkuje symbolicznie. Całka przychodzi jako wartość
  oczekiwana i jest szacowana próbkowaniem, czyli Programem~\ref{prog:F13},
  a potem Programem~\ref{prog:P25} \\
Wielomiany charakterystyczne, Cayley--Hamilton, postać Jordana & Żadna
  biblioteka nie liczy wartości własnej z wielomianu. Program~\ref{prog:P10}
  robi wartości własne jako kierunki niezmiennicze, a Program~\ref{prog:P11}
  rozkład, który naprawdę wywołujesz \\
Wyznaczniki przez rozwinięcie Laplace'a, wzory Cramera & Kosztowne silniowo i
  nigdy nieużywane przy rozmiarze. Program~\ref{prog:P09} uczy, co wyznacznik
  \emph{znaczy}, podaje wzór dla $2 \times 2$ i mówi, że powyżej maszyna liczy
  go eliminacją i ty też powinieneś \\
Częstościowa teoria estymatorów: dostateczność, nierówność Cramér--Rao, UMVU &
  Program~\ref{prog:P26} robi największą wiarygodność jako \emph{cel}, czyli
  tę część, która pojawia się w każdej funkcji straty; teoria optymalności nie
  dochodzi do biurka \\
Klasyczny zwierzyniec testów: tablice $t$, $\chi^{2}$, ANOVA, $F$ & Masz
  komputer. Program~\ref{prog:P27} odpowiada na te same pytania przez losowanie
  ze zwracaniem i przez test dokładny na parach niezgodnych, przy mniejszej
  liczbie założeń i bez tablicy do sprawdzania \\
\bottomrule
\end{tabularx}
\end{center}

\medskip
\textbf{I cztery wyłączenia, co do których ta książka jest najmniej pewna},
wypisane dlatego, że powiedzenie tego teraz jest tańsze niż bronienie ich
później.

\emph{Teoria optymalizacji wypukłej} jest poza zakresem na tej podstawie, że
uczenie głębokie nie jest wypukłe, a kontrargument jest dobry: kwantyzacja,
kalibracja, problemy serwowania i routingu oraz dekodowanie z ograniczeniami są
wypukłe, a inżynier, który nie umie takiego rozpoznać, rozwiązuje go pętlą.
Programy~\ref{prog:P19} i~\ref{prog:P22} zachowują umiejętność rozpoznawania
właśnie dlatego i to jest pierwsze wyłączenie do rewizji, jeśli czytelnicy
zgłoszą, że potrzebują więcej.

\emph{Matematyka dyskretna} zatrzymuje się daleko przed algorytmami i
złożonością. Da się tego bronić tylko dlatego, że czytelnik jest inżynierem
oprogramowania, który spotkał już graf \dash{} czyli przy założeniu, którego ta
książka nie robi nigdzie indziej.

\emph{Wnioskowanie przyczynowe} jest wyłączone całkowicie i jest to coraz
bardziej błędne, w miarę jak ewaluacja i analiza testów A/B stają się częścią
tej roboty. To kandydat na drugie wydanie, a nie temat, który tu nie należy.

\emph{Stochastyczne równania różniczkowe} są wyłączone w tabeli powyżej. Dziś
da się tej granicy bronić i za kilka lat może się nie dać.

\section{Tematy, których uczy, ale tylko do głębokości roboczej}
\label{sec:E-depth}
\index{further reading!working depth}

\begin{warning}
\textbf{To są nazwy, a nazwa jest tu werdyktem tylko tam, gdzie
Sekcja~\ref{sec:E-survey} już go wydała.} Tamta sekcja mówi, co każda praca
robi dobrze, a co źle, bo tamte prace zostały przeczytane i porównane ze sobą,
zanim ta książka powstała. Tam, gdzie nazwa z poniższej listy pojawia się i
tam, jej werdykt zostaje, a ta sekcja nic do niego nie dodaje. Reszta nie
została przeczytana i porównana, a wymyślanie im charakterystyki jest dokładnie
tym ruchem, o którego odrzucenie ta książka prosi przez czterdzieści siedem
programów \dash{} więc dla nich dalej jest wskazówka i nic więcej.
\end{warning}

\medskip
Na \textbf{algebrę liniową} poza roboczą głębokość części~III: Strang, a na
ujęcie abstrakcyjne Axler. Na ten sam teren z nachyleniem ku uczeniu
maszynowemu \dash{} Deisenroth, Faisal i Ong.

Na \textbf{rachunek prawdopodobieństwa} \dash{} Blitzstein i Hwang. Na
\textbf{wnioskowanie} \dash{} Wasserman, \emph{All of Statistics}. Na
\textbf{metody bayesowskie} poza Program~\ref{prog:P28} \dash{} Gelman i
współautorzy, \emph{Bayesian Data Analysis}.

Na \textbf{teorię informacji} poza część~VIII \dash{} Cover i Thomas. Na
\textbf{optymalizację numeryczną} jako przedmiot sam w sobie, a nie jako reguły
aktualizacji, które wyprowadza Program~\ref{prog:P20} \dash{} Nocedal i Wright.

A na uczenie maszynowe, które to wszystko konsumuje: Goodfellow, Bengio i
Courville na obraz z 2016 roku, Murphy i Bishop na obecny.

\medskip
I jedna luka jest tu warta nazwania, zamiast zostawiania jej do odkrycia.
Część~II schodzi o rząd wielkości głębiej niż jedyny rozdział, jaki ten
przegląd znalazł na ten temat, a na samą arytmetykę zmiennoprzecinkową nie
przejrzano tu żadnego źródła, więc ta sekcja nie ma na nią wskazówki.
Sekcja~\ref{sec:E-not} kieruje do Trefethena i Baua
oraz do Highama na numeryczną algebrę liniową, czyli na temat sąsiedni,
a nie ten.

\section{Dokąd prowadzi każde niedowiedzione zdanie}
\label{sec:E-route}
\index{further reading!rigour boxes}
\index{rigour boxes!destinations}

Ramka rygoru nazywa wynik, mówi, że nie jest tu dowodzony, i mówi, gdzie
mieszka dowód. Większość wskazuje inny program tej książki. Te, które wskazują
\emph{poza} nią, nazywają rodzaj kursu \dash{} \enquote{dowolny pierwszy kurs
analizy} \dash{} i, z jednym wyjątkiem, żadnego tytułu. Ta sekcja jest dla nich
skorowidzem i trzymana jest reguła, a nie liczba: ramka, która wychodzi poza
książkę, ma tu wiersz.

\begin{note}
\textbf{Celem jest tutaj praca, którą ten dodatek już przedstawił, albo nic
\dash{} i kolumna mówi, co z tych dwojga.} Tam, gdzie pisze \emph{nie zbadane},
żadne źródło z tego dodatku nie zostało przeczytane pod tym kątem, a nazwanie
jakiegoś mimo to byłoby ruchem, o którego odrzucenie prosi czterdzieści siedem
programów. Ostatnia z trzech tabel poniżej to właśnie te wiersze. To jest stan
zapisu, a nie luka w składzie, i warto to wiedzieć, zanim ruszysz szukać.
\end{note}

\medskip
Najpierw fundamenty pod matematyką, której ta książka uczy.

\begin{center}
\begin{tabularx}{\linewidth}{@{}Xll@{}}
\toprule
\textbf{Powiedziane i niedowiedzione} & \textbf{Odłożone z}
  & \textbf{Dokąd iść} \\
\midrule
Kiedy granica istnieje i kiedy pokrycie nie ma znaczenia &
  \ref{prog:F04}, \ref{prog:F13}, \ref{prog:P09} & Abbott; Tao \\
Różniczkowalność i wzór Taylora z resztą &
  \ref{prog:P15}, \ref{prog:P17} & Abbott; Tao \\
Że dwie bazy mają ten sam rozmiar, i twierdzenie spektralne &
  \ref{prog:P04}, \ref{prog:P10} & Axler; Strang \\
Czym jest prawdopodobieństwo na kontinuum &
  \ref{prog:P23} & Williams; Durrett \\
Że rzeczy niezależne się faktoryzują, i funkcje charakterystyczne &
  \ref{prog:P25} & Blitzstein i Hwang \\
\bottomrule
\end{tabularx}
\end{center}

\medskip
Potem wyniki, których ta książka używa i których nie wyprowadza.

\begin{center}
\begin{tabularx}{\linewidth}{@{}Xll@{}}
\toprule
\textbf{Powiedziane i niedowiedzione} & \textbf{Odłożone z}
  & \textbf{Dokąd iść} \\
\midrule
Warunek regularności ograniczeń i dostateczność &
  \ref{prog:P22} & Boyd i Vandenberghe \\
Nierówność Krafta i granica kodowania źródła &
  \ref{prog:P29}, \ref{prog:P30} & Cover i Thomas \\
Dolna granica dowodowa i wnioskowanie wariacyjne &
  \ref{prog:P19} & Murphy \\
Efektywność asymptotyczna i nierówność Cramér--Rao &
  \ref{prog:P26} & wyłączone powyżej \\
Co ścieżka rezydualna i warstwa normalizująca robią z wariancją &
  \ref{prog:P25} & Program~\ref{prog:P32} \\
\bottomrule
\end{tabularx}
\end{center}

\medskip
I pięć, których ten dodatek nie umie spłacić.

\begin{center}
\begin{tabularx}{\linewidth}{@{}Xll@{}}
\toprule
\textbf{Powiedziane i niedowiedzione} & \textbf{Odłożone z}
  & \textbf{Dokąd iść} \\
\midrule
Metoda gwiazdek i przegród, i zliczanie z powtórzeniami &
  \ref{prog:P12} & nie zbadane \\
Ile naprawdę prawie prostopadłych kierunków mieści przestrzeń &
  \ref{prog:P05} & nie zbadane \\
Dlaczego losowanie ze zwracaniem szacuje rozkład z próby &
  \ref{prog:P27} & nie zbadane \\
Jak napisać dowód, a nie go przeczytać &
  \ref{prog:P14} & nie zbadane \\
Przypadek pierwiastka podwójnego w rekurencji liniowej &
  \ref{prog:P20} & nie zbadane \\
\bottomrule
\end{tabularx}
\end{center}

\medskip
Dwa wiersze tej ostatniej tabeli mają wskazówkę w samym programie, a nie
tutaj: Program~\ref{prog:P27} nazywa pracę, z której wziął się bootstrap, a
Program~\ref{prog:P05} nazywa dziedzinę, do której należy jego pytanie.
Pozostałe nazywają rodzaj kursu i żadnego tytułu. \textbf{Jeden z nich nie jest
pominięciem, tylko własnym stanowiskiem tej książki}: uczy czytać twierdzenie i
mówi wprost, że pisać go nie będziesz musiał, więc podręcznik pisania dowodów
jest jedynym celem, do którego odmawia cię odesłać.

\medskip
To ostatnie zdanie jest sensem całego dodatku. Ta książka nie uczy żadnego z
tych tematów tak głęboko jak jej własne źródła i nigdy nie miała. To, co
twierdzi, jest węższe i sprawdzalne: że możesz je teraz otworzyć.
